\chapter{Couverture de queue et volatilité longue}\label{ch:s2:tail-hedging-and-long-volatility}

Une \omterm{def:s2:tail-hedging-and-long-volatility:hedge}{couverture de queue} qui perd un peu d'argent chaque mois est pendant des années la position la plus détestée d'un portefeuille, jusqu'au mois où elle les rémunère toutes. L'indice de Cboe
5~\% Put Protection Index, qui détient le S\&P 500 et achète chaque mois un put de 5~\% hors de la monnaie, perd 10.6~\% en octobre 1987 alors que l'indice chute
de 21.8~\%, et 0.4~\% sur février et mars 2020, lorsque l'indice recule de 19.9~\%. De juillet 2007 à mars 2009, il perd tout de même 42.0~\%, car les puts mensuels
protègent contre un krach, mais non contre un marché baissier lent. Sur le marché synthétique de ce chapitre, un programme de puts mensuels coûte 6.0~\% par an les années
sans krach, rapporte 19.2~\% et 39.5~\% dans deux des trois années de krach, et fait croître le portefeuille de 2.5~\% par an contre 3.0~\% sans couverture,
avec un drawdown maximal ramené de 73~\% à 41~\%. Le fait que cela en valait la peine dépend de ce que la couverture permet au reste du portefeuille de faire. La construction est
\texttt{firm.tailhedge}.

\section{Le coût de la convexité}

\begin{definition}[Couverture de queue]\label{def:s2:tail-hedging-and-long-volatility:hedge}
Une \emph{couverture de queue}\index{couverture de queue} (tail hedge) est une position conçue pour rapporter lors de chutes importantes et rares du marché, comme des puts hors de la monnaie, des put spreads, une variance longue
ou des contrats à terme sur la volatilité, achetée à un coût en temps ordinaire en échange d'une convexité lors d'un krach.
\end{definition}

\begin{definition}[Coût de portage de la couverture]\label{def:s2:tail-hedging-and-long-volatility:bleed}
Le \emph{coût de portage de la couverture}\index{coût de portage de la couverture} (bleed) est le coût d'une \omterm{def:s2:tail-hedging-and-long-volatility:hedge}{couverture de queue} sur les périodes sans krach~: les primes payées moins les gains perçus, exprimés en part du
portefeuille, par an.
\end{definition}

Un put hors de la monnaie se price à son point sur le smile, qui pour les puts sur indice se situe au-dessus de la volatilité à la monnaie. Cette volatilité intègre déjà la prime
de variance du chapitre~1. L'acheteur paie donc les deux primes, celle que le vendeur du chapitre~1 encaisse ainsi que la prime elle-même. Harvey et ses coauteurs ont comparé des stratégies défensives lors des
grandes baisses boursières et des récessions. Détenir en continu des puts à courte échéance sur le S\&P 500 a été la défense la plus fiable, et aussi la plus coûteuse.

\texttt{firm.tailhedge} fait tourner des programmes de puts sur l'indice synthétique du chapitre~1. Il couvre vingt ans, avec trois krachs programmés de 20~\% à 30~\%
en dix jours, aux années 5, 12 et 17. Les options sont pricées sur un smile plus pentu que celui du chapitre~1 (skew $-0.25$ par unité standardisée, courbure 0.02). Le
portefeuille détient l'indice et un put par unité d'indice, et à chaque roll il fixe ses unités de sorte que l'action et les puts couvrent le même notionnel
(\cref{lst:s2:tail-hedging-and-long-volatility:programme}).

\section{Programmes de puts et leur monétisation}

\begin{definition}[Monétisation de la couverture]\label{def:s2:tail-hedging-and-long-volatility:monetisation}
La \emph{monétisation de la couverture}\index{monétisation de la couverture} (hedge monetisation) est la règle qui consiste à vendre une \omterm{def:s2:tail-hedging-and-long-volatility:hedge}{couverture de queue} ayant pris de la valeur, avant son échéance, puis à la remplacer au
nouveau niveau de marché~: par exemple, lorsque le put vaut trois fois son coût d'achat. Elle transforme un gain latent en liquidités réinvestissables, et
rétablit la protection près du nouveau spot.
\end{definition}

Les indices de Cboe montrent la forme réelle (\cref{fig:s2:tail-hedging-and-long-volatility:episodes}). De juin 1986 à septembre 2026, PPUT, qui est un
indice de rendement total, a progressé au taux logarithmique de 7.7~\% par an, avec une volatilité de 13.5~\%. CLLZ détient le S\&P 500, achète chaque mois un put spread de 2.5~\%--5~\%
et le finance en vendant un call. Il a progressé de 7.8~\%, avec une volatilité de 14.6~\%. L'indice de prix du S\&P 500 a progressé de 8.5~\%, avec une volatilité de 18.4~\%.
L'indice de prix omet les dividendes, si bien que l'écart de rendement sous-estime le coût des couvertures. Les 21 pires jours de PPUT ont coûté 13.5~\%, contre 33.0~\% pour l'indice. Lors du
krach de l'automne 2008, de septembre à novembre, PPUT a perdu 6.7~\%, l'indice 30.1~\% et CLLZ 31.2~\%~: un put spread ne protège que la baisse comprise entre ses
strikes. Les plus grands drawdowns en disent davantage. PPUT a perdu 42.0~\% de juillet 2007 à mars 2009, quand l'indice en perdait 56.8~\%. La chute a duré vingt
mois, et un put dont le strike se situe 5~\% sous le début de chaque mois ne couvre que la part de la baisse de ce mois qui excède 5~\%.

\begin{omfigure}
\begin{tikzpicture}
\begin{axis}[width=11cm,height=5.6cm,ybar,bar width=8pt,ylabel={\small rendement (\%)},tick label style={font=\footnotesize},
  xmin=0.5,xmax=4.5,ymin=-60,ymax=5,xtick={1,2,3,4},xticklabels={oct. 1987,automne 2008,début 2020,drawdown maximal},
  grid=major,grid style={gray!25},table/col sep=comma,legend style={font=\scriptsize,at={(0.5,-0.2)},anchor=north},legend columns=3]
\addplot[fill=omThm!70,draw=omThm] table[x=x,y=PPUT]{figdata/strategies-2/07-tail-hedging-and-long-volatility/episodes.csv};
\addplot[fill=omDef!60,draw=omDef] table[x=x,y=CLLZ]{figdata/strategies-2/07-tail-hedging-and-long-volatility/episodes.csv};
\addplot[fill=gray!40,draw=gray] table[x=x,y=SPX]{figdata/strategies-2/07-tail-hedging-and-long-volatility/episodes.csv};
\legend{PPUT (puts mensuels de 5~\%),CLLZ (collar à put spread),indice de prix S\&P 500}
\end{axis}
\end{tikzpicture}

\omcaption{Indices de protection par puts et de collar de Cboe comparés au S\&P 500 lors de trois krachs et sur le plus grand drawdown de chaque série (tous trois de
2007 à mars 2009), 1986--2026. Les puts mensuels protègent contre les baisses rapides, non les baisses lentes. Dérivé des historiques d'indices de Cboe~; les séries brutes ne sont pas
redistribuées. Données~: \texttt{s2\_fetch\_protect}.}\label{fig:s2:tail-hedging-and-long-volatility:episodes}
\end{omfigure}

Sur le marché synthétique, les programmes se comparent comme suit, chacun face à l'indice et face à l'alternative la plus simple, une allocation statique de 70~\% en indice
et 30~\% en cash, rééquilibrée chaque mois~:

\begin{center}
\small
\begin{tabular}{@{}p{4.6cm}>{\centering\arraybackslash}p{2.5cm}>{\centering\arraybackslash}p{2.5cm}>{\centering\arraybackslash}p{2.5cm}>{\centering\arraybackslash}p{2.5cm}@{}}
\toprule
vingt ans & croissance (\%/an) & volatilité (\%) & drawdown maximal (\%) & pire mois (\%)\\
\midrule
indice & 3.0 & 19.0 & $-73$ & $-43$\\
puts 5~\%, mensuels & 2.5 & 13.9 & $-41$ & $-17$\\
puts 10~\%, mensuels & 3.2 & 16.0 & $-49$ & $-22$\\
puts 10~\%, trimestriels & 3.0 & 14.0 & $-46$ & $-24$\\
puts 20~\%, trimestriels & 4.2 & 16.2 & $-52$ & $-33$\\
10~\% trimestriels, monétisés à 3$\times$ & 2.2 & 14.0 & $-62$ & $-26$\\
puts 20~\%, semestriels & 4.6 & 15.5 & $-50$ & $-36$\\
70~\% indice, 30~\% cash & 2.8 & 13.2 & $-58$ & $-31$\\
indice + overlay de tendance & 6.3 & 17.6 & $-56$ & $-24$\\
\bottomrule
\end{tabular}
\end{center}

La comptabilité annuelle est celle de la couverture elle-même. Sur les 17 années calmes, le programme de puts mensuels à 5~\% a coûté 6.0~\% du portefeuille par an. Lors des trois
années de krach, il a rapporté 19.2~\%, perdu 2.7~\% puis gagné 39.5~\%. Le second krach est survenu alors que la volatilité de l'indice valait 3.9~\% et que le put était bon marché, mais
le put n'a rapporté que 12.0~\% du portefeuille pendant les semaines mêmes du krach, contre 30.0~\% et 29.2~\% dans les deux autres. Après un krach, la volatilité reste élevée pendant des mois, et les puts achetés alors sont chers~: le coût de portage de cette année-là a absorbé le
gain. Des puts plus éloignés coûtent moins cher à porter (1.5~\% par an pour des puts semestriels à 20~\%) et paient plus tard et moins (2.0~\%, $-2.1~\%$ et 26.3~\%). Le programme monétisé
a vendu ses puts onze fois à trois fois leur coût. Il a rapporté le plus lors de la première année de krach (16.3~\%), en vendant tôt à l'automne. Il n'a ensuite conservé qu'une protection dont le strike se situait près
du nouveau spot, plus bas, et le reste de la baisse est passé en dessous~: son drawdown, 62~\%, est le pire des programmes.

\section{Alternatives~: tendance, volatilité longue, variance}

Harvey et ses coauteurs ont constaté que le momentum temporel (time-series momentum) sur contrats à terme a bien résisté lors des longues phases de vente, et qu'une stratégie long--short de qualité a bien résisté lors des
fuites vers la qualité. Les bons du Trésor portent un carry positif, mais la corrélation négative entre obligations et actions observée après 2000 reste une rareté historique. L'or en position longue et la
protection de crédit longue se situaient entre les puts et les obligations, en coût comme en fiabilité. Sur l'indice synthétique, un overlay de tendance ajoute un notionnel égal à la moitié de celui de l'indice, dans le
sens du rendement de l'année écoulée. Il croît de 6.3~\% par an, avec un drawdown de 56~\% et un pire mois de $-24~\%$~: meilleur que les puts en croissance,
moins bon en drawdown. Il lui faut une phase de vente assez longue pour que le signal se retourne. L'alpha de crise de ce type (livre~8, chapitre~19) et
la convexité d'une option ne couvrent pas les mêmes risques.

La variance longue et les contrats à terme longs sur le VIX sont les versions, propres au livre de volatilité, de la couverture. Ils rapportent la prime de variance et le roll-down du chapitre~5 en
temps ordinaire, et versent un gain lors d'un pic de volatilité, que l'indice fasse ou non un gap. Ils requièrent leur propre règle de monétisation, car une position en variance
peut voir son gain disparaître en quelques semaines dès que la volatilité retombe.

\section{Juger une couverture à l'aune du portefeuille}

L'avertissement d'Israelov est le test à appliquer. Dans l'usage courant, les puts se sont montrés plutôt inefficaces pour réduire les drawdowns, en comparaison de l'alternative simple
qui consiste à détenir statiquement une part moindre de l'indice. À moins que les achats et les échéances ne soient calés avec soin autour du drawdown, ils peuvent aggraver les drawdowns et la
volatilité par unité de rendement espéré. Le marché synthétique est clément envers les puts~: ses krachs sont rapides, dix jours chacun, ce pour quoi un put est
conçu. Là, le programme mensuel de 5~\% bat le mélange statique sur le drawdown ($-41~\%$ contre $-58~\%$ à volatilité voisine) et perd sur la croissance (2.5~\% contre
2.8~\%). Le pire drawdown de l'indice réel a été lent, et PPUT en a perdu 42.0~\%.

Une couverture se juge donc à ce qu'elle permet au portefeuille de faire. Avec une couverture qui plafonne la perte, le reste du portefeuille peut détenir davantage de risque, si bien que le coût
revient sous la forme d'une exposition supplémentaire. Sans cela, la couverture n'est qu'un frein habillé d'un bon argument. Le taux de croissance est la mesure retenue, sur un historique assez long pour
contenir des krachs des deux types.

\begin{omfigure}
\begin{tikzpicture}
\begin{axis}[width=11cm,height=5.6cm,ymode=log,xlabel={\small année},ylabel={\small valeur (départ 1)},tick label style={font=\footnotesize},
  xmin=0,xmax=20,ymin=0.7,ymax=5,ytick={0.75,1,1.5,2,3,4},log ticks with fixed point,grid=major,grid style={gray!25},table/col sep=comma,
  legend style={font=\scriptsize,at={(0.5,-0.3)},anchor=north},legend columns=4]
\addplot[gray,thick] table[x=year,y=index]{figdata/strategies-2/07-tail-hedging-and-long-volatility/paths.csv};
\addplot[omThm,thick] table[x=year,y=puts]{figdata/strategies-2/07-tail-hedging-and-long-volatility/paths.csv};
\addplot[omDef,thick,dashed] table[x=year,y=static]{figdata/strategies-2/07-tail-hedging-and-long-volatility/paths.csv};
\addplot[black,thin] table[x=year,y=trend]{figdata/strategies-2/07-tail-hedging-and-long-volatility/paths.csv};
\legend{indice,puts mensuels à 5~\%,70~\% indice + cash,indice + tendance}
\end{axis}
\end{tikzpicture}

\omcaption{L'indice synthétique et trois portefeuilles protégés sur vingt ans, avec des krachs aux années 5, 12 et 17, en échelle logarithmique~: un programme mensuel de
puts à 5~\%, une part statique de 70~\% dans l'indice, et l'indice assorti d'une surcouche de tendance. Données~: \texttt{s2\_tailhedge.curves}.}\label{fig:s2:tail-hedging-and-long-volatility:paths}
\end{omfigure}

\section{Fiches de stratégie}

\begin{strategyfile}{Programme glissant de puts hors de la monnaie}\label{strat:s2:tail-hedging-and-long-volatility:puts}
\sfield{Qui paie, et pourquoi} Personne, en moyenne~: l'acheteur paie les primes de variance et de skew en échange de la convexité lors d'un krach.
\sfield{Instruments et plateformes} Puts indiciels cotés, d'échéance mensuelle à semestrielle, à 5~\% à 20~\% hors de la monnaie.
\sfield{Signal} Aucun~; c'est un programme, non une transaction. Certains programmes achètent davantage lorsque la protection est bon marché.
\sfield{Dimensionnement et exécution} Un budget annuel fixe, reconduit mécaniquement~; les règles de monétisation sont décidées à l'avance.
\sfield{Coûts} Primes et fourchettes~; le \omterm{def:s2:tail-hedging-and-long-volatility:bleed}{coût de portage de la couverture}.
\sfield{Comment il meurt} Les marchés baissiers lents qui restent au-dessus du strike de chaque put~; l'abandon après des années de coût de portage, juste avant qu'il paie.
\sfield{Horizon, capacité, infrastructure} Des années~; grande capacité sur les options sur indices.
\sfield{Backtest honnête} Inclure les drawdowns lents comme rapides~; comparer avec la détention statique d'une part moindre.
\sfield{Sources} Israelov (2019)~; Harvey et coauteurs (2019)~; Cboe PPUT~: $-42.0~\%$ de juillet 2007 à mars 2009~; ce chapitre~: un coût de portage de 6.0~\% par an
et un drawdown ramené de 73~\% à 41~\%.
\end{strategyfile}

\begin{strategyfile}{Collar à spread de puts}\label{strat:s2:tail-hedging-and-long-volatility:collar}
\sfield{Qui paie, et pourquoi} L'acheteur de call finance le spread de puts~; le détenteur renonce à la hausse pour plafonner une bande de baisse.
\sfield{Instruments et plateformes} Spreads de puts et calls indiciels, mensuels.
\sfield{Signal} Aucun~; un programme.
\sfield{Dimensionnement et exécution} Coût nul par construction~: le strike du call est fixé de sorte que sa prime finance le spread de puts.
\sfield{Coûts} Fourchettes sur trois options par mois~; la hausse abandonnée.
\sfield{Comment il meurt} Un krach au-delà du strike inférieur du spread de puts~: CLLZ a perdu 31.2~\% à l'automne 2008.
\sfield{Horizon, capacité, infrastructure} Des mois~; un roll mécanique.
\sfield{Backtest honnête} Mesurer les krachs sous le strike inférieur, et non seulement la bande.
\sfield{Sources} Cboe CLLZ (statistiques dérivées dans ce chapitre).
\end{strategyfile}

\begin{strategyfile}{Variance longue avec monétisation}\label{strat:s2:tail-hedging-and-long-volatility:variance}
\sfield{Qui paie, et pourquoi} Personne, en moyenne~: elle paie la prime de variance en échange d'un gain lors des pics de volatilité.
\sfield{Instruments et plateformes} Swaps de variance, contrats à terme sur le VIX ou calls.
\sfield{Signal} Volatilité bon marché par rapport à son historique~; la pente de la courbe du VIX.
\sfield{Dimensionnement et exécution} Un budget~; monétiser à un multiple fixé à l'avance, ou lorsque la volatilité a bondi et que la courbe s'inverse.
\sfield{Coûts} La prime de variance~; le roll-down des contrats à terme sur le VIX.
\sfield{Comment il meurt} Les pics qui retombent avant la monétisation~; un krach sans grande volatilité.
\sfield{Horizon, capacité, infrastructure} Des mois~; accès à la variance ou au VIX.
\sfield{Backtest honnête} Utiliser le cours au moment de la monétisation, et non le sommet.
\sfield{Sources} Harvey et coauteurs (2019) sur les stratégies défensives~; aucune performance chiffrée vérifiée.
\end{strategyfile}

\section{Tutoriel~: payer pour le mois qui paie}\label{tut:s2:tail-hedging-and-long-volatility:tut}

\begin{tutorial}
\textbf{Objectif.} Faire tourner des programmes de puts à plusieurs strikes, échéances et règles de monétisation sur l'indice synthétique, les comparer avec un mélange statique et une
surcouche de tendance, et lire les indices de protection de Cboe. \textbf{État final~:} le tableau et les deux figures.
\begin{tutsteps}
\item \textbf{Le programme}.
\omcode{code/firm/tailhedge/firm_tailhedge.py}{43}{63}{Un programme glissant de puts avec monétisation, entièrement investi.}\label{lst:s2:tail-hedging-and-long-volatility:programme}
\item \textbf{Le coût de portage}.
\omcode{code/strategies-2/07-tail-hedging-and-long-volatility/python/s2_tailhedge.py}{57}{72}{Résultat annuel de la couverture de chaque programme, années calmes et années de krach.}\label{lst:s2:tail-hedging-and-long-volatility:bleed}
\item \textbf{Exécuter} \texttt{s2\_fetch\_protect.py} une fois, puis \texttt{table()}, \texttt{bleed()} et \texttt{fig\_tailhedge.py}.
\end{tutsteps}
\textbf{Pour aller plus loin.} Introduire un marché baissier lent (une année de baisse régulière) et relancer le calcul~; mettre à levier le portefeuille couvert à la volatilité de celui qui ne l'est pas,
comparer la croissance~; n'acheter des puts que lorsque le skew est plat.
\end{tutorial}

\section{Construction~: couverture de queue}\label{bld:s2:tail-hedging-and-long-volatility:tailhedge}

\begin{build}
\bfield{Rôle} Programmes glissants de puts avec monétisation, mélanges statiques et surcouches de tendance, jugés sur la croissance, la volatilité et le drawdown.
\bfield{Interface} \texttt{put\_vol(S, K, tau, v, cfg)}, \texttt{put\_programme(r, v, cfg, otm, tenor, monetise)},
\texttt{static\_mix(r, weight, every)}, \texttt{trend\_overlay(r, weight, lookback)}, \texttt{evaluate(nav)}.
\bfield{Règles} Puts évalués et valorisés sur la surface, avec le smile tronqué~; entièrement investi à chaque roll~; la monétisation remplace le put au
nouveau spot.
\bfield{Tests d'acceptation} \texttt{code/firm/tailhedge/tests/}~: un marché immobile ne fait que coûter~; le put couvre une baisse sous son strike~; un mélange statique, une
surcouche de tendance et un drawdown calculé à la main.
\bfield{Pour aller plus loin} Spreads de puts et collars~; couvertures en variance et sur le VIX~; levier vers une volatilité cible.
\end{build}

\begin{omsources}
\item R.~Israelov, ``Pathetic protection: the elusive benefits of protective puts'', \emph{Journal of Alternative Investments} 21(3), 2019.
\item C.~R.~Harvey, E.~Hoyle, S.~Rattray, M.~Sargaison, D.~Taylor and O.~Van Hemert, ``The best of strategies for the worst of times: can portfolios be
crisis proofed?'', \emph{Journal of Portfolio Management} 45(5), 2019.
\item Cboe, S\&P 500 Put Protection Indices, méthodologie~; Cboe Insights sur PPUT, CLL et CLLZ, 2021.
\item Cboe Global Markets, historiques quotidiens de PPUT, CLLZ et SPX.
\end{omsources}

\section{Exercices}

\begin{exercise}[$\star$]\label{exo:s2:tail-hedging-and-long-volatility:1}
Un programme de puts coûte 6.0~\% par an dans les années calmes. Sur 17 années calmes et trois années de krach rapportant 19.2~\%, $-2.7~\%$ et 39.5~\%.
Calculer son total, sans tenir compte de la capitalisation.
\end{exercise}

\begin{exercise}[$\star$]\label{exo:s2:tail-hedging-and-long-volatility:2}
Expliquer pourquoi un put mensuel à 5~\% ne protège pas contre une baisse de 1~\% par semaine pendant un an.
\end{exercise}

\begin{exercise}[$\star$]\label{exo:s2:tail-hedging-and-long-volatility:3}
Un put coûte 1~\% du notionnel et est monétisé à trois fois son coût. Déterminer sa valeur à la vente.
\end{exercise}

\begin{exercise}[$\star\star$]\label{exo:s2:tail-hedging-and-long-volatility:4}
Pourquoi CLLZ a-t-il perdu presque autant que l'indice à l'automne 2008 alors que PPUT n'a perdu que 6.7~\%~?
\end{exercise}

\begin{exercise}[$\star\star$]\label{exo:s2:tail-hedging-and-long-volatility:5}
Pourquoi l'écart de rendement entre PPUT et l'indice S\&P 500 de cours sous-estime-t-il le coût de la couverture~?
\end{exercise}

\begin{exercise}[$\star\star$]\label{exo:s2:tail-hedging-and-long-volatility:6}
En quoi la protection apportée par un overlay de tendance diffère-t-elle de celle d'un put~?
\end{exercise}

\begin{exercise}[$\star\star\star$]\label{exo:s2:tail-hedging-and-long-volatility:7}
\emph{Programmation.} Exécuter \texttt{slow\_bear()}, qui ajoute une baisse régulière de 40~\% en log-rendement sur la huitième année de l'indice synthétique, et comparer
la perte de l'année du programme mensuel à 5~\%, du mix statique et de l'indice. Lequel protège le mieux, et pourquoi~?
\end{exercise}

\begin{exercise}[$\star\star\star$]\label{exo:s2:tail-hedging-and-long-volatility:8}
\emph{Trouver la faille.} «~La couverture nous a coûté 6~\% par an pendant cinq ans~; elle ne fonctionne pas, on la supprime.~»
\end{exercise}

\section{Problème~: payer pour le mois qui rapporte}

\begin{problem}[{Problème du week-end --- un programme de couverture de queue}]\label{pb:s2:tail-hedging-and-long-volatility:1}
Le marché synthétique du chapitre, les indices de protection de Cboe et le bilan public.

\medskip\noindent\textbf{Partie I --- Le coût.}
\begin{enumerate}
\item Définir une \omterm{def:s2:tail-hedging-and-long-volatility:hedge}{couverture de queue} et son coût de portage.
\item Quelles primes paie l'acheteur d'un put~?
\item Que Harvey et ses coauteurs ont-ils établi au sujet des puts et des autres défenses~?
\item Qu'a établi Israelov~?
\end{enumerate}

\medskip\noindent\textbf{Partie II --- Les indices réels.}
\begin{enumerate}[resume]
\item Donner la croissance et la volatilité de PPUT, de CLLZ et du S\&P 500.
\item Comment chacun s'est-il comporté en 1987, en 2008 et en 2020~?
\item Pourquoi PPUT a-t-il perdu 42~\% entre 2007 et 2009~?
\item Pourquoi un put spread échoue-t-il dans un krach profond~?
\end{enumerate}

\medskip\noindent\textbf{Partie III --- Les programmes.}
\begin{enumerate}[resume]
\item Définir la monétisation.
\item Comparer la croissance et les drawdowns des programmes.
\item Donner le coût de portage et les gains des années de krach du programme mensuel à 5~\%.
\item Pourquoi la monétisation a-t-elle aggravé le drawdown~?
\end{enumerate}

\medskip\noindent\textbf{Partie IV --- Le verdict.}
\begin{enumerate}[resume]
\item Énoncer le \emph{résultat nommé}~: le coût annuel de la couverture et le taux de croissance du portefeuille avec et sans elle.
\item Comment la couverture se compare-t-elle à un mix statique~?
\item Comment se compare-t-elle à un overlay de tendance~?
\item Pourquoi le marché synthétique est-il favorable aux puts~?
\item Comment faut-il juger une couverture~?
\item Comment backtester honnêtement une \omterm{def:s2:tail-hedging-and-long-volatility:hedge}{couverture de queue}~?
\item Quelle stratégie du dossier convient le mieux à un marché baissier lent~?
\item En une phrase~: que achète une \omterm{def:s2:tail-hedging-and-long-volatility:hedge}{couverture de queue}~?
\end{enumerate}
\end{problem}

\section{Questions d'entretien}

\begin{interviewq}[$\star$ \iqroles{trader}]\label{iq:s2:tail-hedging-and-long-volatility:1}
Pourquoi les puts sur indice hors de la monnaie sont-ils chers~?
\end{interviewq}

\begin{interviewq}[$\star\star$ \iqroles{chercheur}]\label{iq:s2:tail-hedging-and-long-volatility:2}
Comment compareriez-vous un programme de puts à la détention d'une part moindre d'actions~?
\end{interviewq}

\begin{interviewq}[$\star\star$ \iqroles{trader}]\label{iq:s2:tail-hedging-and-long-volatility:3}
Vos puts valent cinq fois ce que vous avez payé. Les vendez-vous~?
\end{interviewq}

\begin{interviewq}[$\star\star$ \iqroles{risque}]\label{iq:s2:tail-hedging-and-long-volatility:4}
Quels scénarios utiliseriez-vous pour tester un programme de \omterm{def:s2:tail-hedging-and-long-volatility:hedge}{couverture de queue}~?
\end{interviewq}

\begin{interviewq}[$\star\star$ \iqroles{développeur}]\label{iq:s2:tail-hedging-and-long-volatility:5}
Concevez le processus quotidien d'un programme de puts glissants avec des règles de monétisation.
\end{interviewq}

\begin{interviewq}[$\star\star\star$ \iqroles{chercheur}]\label{iq:s2:tail-hedging-and-long-volatility:6}
Montrez que, pour un portefeuille à rendements log-normaux, le taux de croissance vaut approximativement $\mu - \sigma^2/2$, et expliquez comment une couverture qui abaisse $\sigma$ peut
relever le taux de croissance même si elle abaisse $\mu$.
\end{interviewq}
